
%  APPENDIX MATERIAL BELONGING TO SEC. 6.  Everything below this line is NOT
%  part of the body section: move it verbatim into the appendix file, where it
%  must carry \label{app:witness} (that label already exists in theorem_b2.tex
%  and is reused here for the same content).  Sec. 6 references it twice.
% =============================================================================
\section{The geminal witness: proof and angle sweep}
\label{app:witness}

\begin{proof}[Proof of Theorem~\ref{thm:witness}]
Each geminal occupies its own four spin-orbitals, so \eqref{eq:apsg} is a product across the $K$ blocks
and $\gamma$ is block diagonal with spatial natural occupations
$(n_{b_i},n_{a_i})=(2\cos^2\theta,2\sin^2\theta)$; hence $N_{\mathrm u}=8K\cos^2\theta\sin^2\theta=2K$ at
$\theta=\pi/4$, and $\mathcal F_1=4K$ by Eq.~\eqref{eq:collapse}. The Majorana covariance is likewise
block diagonal, and at $\theta=\pi/4$ a single maximally-paired geminal is the two-determinant cat, whose
covariance vanishes identically ($\Gamma=0$): every order saturates its maximum,
$\mathcal F_k=2n_{\rm o}=4$ per geminal for all $k$, giving $\mathcal F_2=4K$. This equality of orders is
a knife-edge of maximal pairing; away from it the orders separate ($\theta=\pi/6$ gives
$\mathcal F_1=3$, $\mathcal F_2=3.75$ per pair, both still large on the same family). In the pairing basis
each geminal is a two-configuration superposition, so the product carries $2^{K}$ determinants. Finally,
as a product of two-orbital factors the state is a matrix-product state whose every bond crosses at most
one geminal, so $\chi=2$ in that ordering, and contraction and Born sampling then cost $O(K)$. All four
values are reproduced by exact diagonalization for $K=1,\dots,4$.
\end{proof}

Table~\ref{tab:fibersweep} continues the fiber exhibit of Table~\ref{tab:fiber} into a continuous sweep of
the rotation angle at $K=3$. It stops at $s=0.654$: at the maximal-pairing angle $s=\pi/4$ the Born
distribution of $\phi$ is exactly flat over $2^{K}$ determinants, so the top-$|\mathcal S|$ rule does not
select a subspace and no contrast measured there is well defined. Across every row the captured weight is
identical at the two endpoints to $2.2\times10^{-16}$, so the entire difference sits in the leakage term,
and the certificate at the rotated endpoint is the trivial bound $1+w_S$ in every row.

\begin{table*}[tp]\centering\small
\caption{\textbf{Angle sweep of the fiber exhibit}, $K=3$, $\delta=0.01$, $\eta=0.10$,
$|\mathcal S|=4$. Columns as in Table~\ref{tab:fiber}.}
\label{tab:fibersweep}
\setlength{\tabcolsep}{4pt}
\begin{tabular}{@{}lccccc@{}}
\toprule
$s$ & $w_S$ (both endpoints)
 & $\widehat\Lambda/\eta$ al. & $\widehat\Lambda/\eta$ rot.
 & cert.\ al. & cert.\ rot. \\
\midrule
$0.131$ & $0.982963$ & $0.328$ & $1.610$ & $\mathbf{0.913}$ & $1.983$ (trivial) \\
$0.262$ & $0.933013$ & $0.317$ & $3.010$ & $\mathbf{1.131}$ & $1.933$ (trivial) \\
$0.393$ & $0.853553$ & $0.299$ & $4.124$ & $\mathbf{1.312}$ & $1.854$ (trivial) \\
$0.524$ & $0.750000$ & $0.277$ & $4.849$ & $\mathbf{1.450}$ & $1.750$ (trivial) \\
$0.654$ & $0.629410$ & $0.252$ & $5.139$ & $\mathbf{1.543}$ & $1.629$ (trivial) \\
\bottomrule
\end{tabular}
\end{table*}
